\uuid{yaqy}
\exo7id{4746}
\titre{exo7 4746}
\auteur{quercia}
\organisation{exo7}
\datecreate{2010-03-16}
\isIndication{false}
\isCorrection{false}
\chapitre{Topologie}
\sousChapitre{Topologie des espaces vectoriels normés}
\cpge{mpsi-mp2i.dimension-finie.a, pcsi.dimension-finie.a, ptsi.dimension-finie.a, mpsi-mp2i.dimension-finie.b, pcsi.dimension-finie.b, ptsi.dimension-finie.b, bcpst1.espace-vectoriel-kn.a, mp-mpi.topologie-espaces-normes.b, pc.espaces-vectoriels-normes.b, psi.espaces-vectoriels-normes.b, mp-mpi.topologie-espaces-normes.j, pc.espaces-vectoriels-normes.g, psi.espaces-vectoriels-normes.g, bcpst2.espaces-vectoriels.a}
\module{Analyse}
\niveau{L2}
\difficulte{}

\contenu{
\question{
Soit $E$ un evn de dimension finie et $({\vec u}_n)$, $({\vec v}_n)$
deux suites de vecteurs telles que :
$$\forall\ n \in \N,\ {\vec u}_n \text{ est colin{\'e}aire {\`a} } {\vec v}_n,
  \qquad
  {\vec u}_n \xrightarrow[n\to\infty]{} \vec u,  \qquad
  {\vec v}_n \xrightarrow[n\to\infty]{} \vec v.$$
Montrer que $\vec u$ et $\vec v$ sont colin{\'e}aires
(raisonner par l'absurde et compl{\'e}ter $(\vec u,\vec v)$ en une base de $E$).
}
}
